\subsection{根基的进一步性质}

命题 5. 设 g 是李代数，r 是其根基。

\begin{enumerate}
\def\labelenumi{(\alph{enumi})}
\item
  如果 \(\rho\) 是 \(\mathfrak{g}\) 的有限维表示，\(\beta\) 是相应的双线性形式，则 \(\mathfrak{r}\) 和 \(\mathcal{D}\mathfrak{g}\) 关于 \(\beta\) 正交。
\item
  r 关于 Killing 形式是 \(\mathcal{D}\mathfrak{g}\) 的正交子空间。
\end{enumerate}

设 \(x, y\) 属于 \(\mathfrak{g}, z \in \mathfrak{r}\)。则 \([y, z] \in \mathcal{D}\mathfrak{g} \cap \mathfrak{r}\)，所以

\[
\beta ([ x, y ], z) = \beta (x, [ y, z ]) = 0
\]

（定理 1）。所以 (a)。

令 \(\mathfrak{r}'\) 是 \(\mathcal{D}\mathfrak{g}\) 关于 Killing 形式的正交子空间。它是 \(\mathfrak{g}\) 的理想（\(\S 3\)，第 6 号，命题 7 (a)），由上面所述包含 \(\mathfrak{r}\)。另一方面，\(\mathfrak{r}'\) 在 \(\mathfrak{g}\) 的伴随表示下的像 s 是可解的（定理 2），所以 \(\mathfrak{r}'\) 是可解的，因为它是 s 的中心扩张。因此 \(\mathfrak{r}' \subset \mathfrak{r}\)。

推论 1. 设 \(\mathfrak{g}\) 是李代数。则 \(\mathfrak{g}\) 可解，当且仅当 \(\mathcal{D}\mathfrak{g}\) 关于 Killing 形式与 \(\mathfrak{g}\) 正交。

这立即由命题 5 (b) 得到。

推论 2. 李代数 \(\mathfrak{r}\) 的根基 \(\mathfrak{g}\) 是特征理想。

\(\mathcal{D}\mathfrak{g}\) 是特征理想，且 Killing 形式是完全不变的（§ 3，第 6 号，命题 10）。因此，\(\mathcal{D}\mathfrak{g}\) 关于 Killing 形式的正交子空间是特征理想（§ 3，第 6 号，命题 7 (b)）。

推论 3。设 g 是李代数，r 是其根基，a 是 g 的理想。则 a 的根基等于 r 与 a 的交。

\(r \cap a\) 是 \(a\) 的可解理想，因而包含于 \(r'\) 的根基 \(a\) 中。反之，\(r'\) 是 \(g\) 的理想（推论 2 以及 §1，第 4 号，命题 2），所以 \(r' \subset r\)。

推论 2 可以精确表述如下：

命题 6. 设 \(\mathfrak{g}\) 是李代数，\(\mathfrak{r}\) 是其根基，\(\mathfrak{n}\) 是其最大幂零理想。每个 \(\mathfrak{g}\) 的导子都将 \(\mathfrak{r}\) 映入 \(\mathfrak{n}\)。

设 D 是 g 的导子。令 \(g' = g + Kx_0\) 是一个李代数，其中 g 是余维 1 的理想，且对所有 \(Dx = [x_0, x]\) 有 \(x \in g\)（§ 1，第 8 号，例 1）。根据命题 5 的推论 3，r 包含于 \(r'\) 的根基 \(g'\) 中。于是 \(D(r) = [x_0, r] \subset [g', g'] \cap r' = s'\)。对所有 \(x \in s' \cup g\)，\(ad_g x\) 都是幂零的（定理 1）。因此，对所有 \(x \in s' \cap g\)，\(ad_g x\) 都是幂零的。所以 \(D(r)\) 包含于 g 的幂零理想 \(s' \cap g\) 中。

推论。李代数的最大幂零理想是特征理想。

注记。总结上述一些结果：若 r、n、s、t 分别表示 g 的根基、g 的最大幂零理想、g 的幂零根基以及 g 关于 Killing 形式的正交子空间，则

\[
\mathfrak {r} \supset \mathfrak {k} \supset \mathfrak {n} \supset \mathfrak {s}.
\]

包含关系 \(\mathfrak{r} \supset \mathfrak{t}\) 由命题 5 (b) 得到。包含关系 \(\mathfrak{t} \supset \mathfrak{n}\) 由 §4，第 4 号，命题 6 (b) 得到。包含关系 \(\mathfrak{n} \supset \mathfrak{s}\) 已在第 3 号的注记 2 中指出。

